\subsection{Prime Ideals}

\begin{definition} \label{def:prime-ideal} \hlt{Prime Ideals}\\
Let $R$ be a commutative ring with identity. An ideal $\mathfrak{p}\trianglelefteq R$ is \hlt{prime} if $\mathfrak{p}\neq R$ and $(\forall\,a,b\in R)\,[ab\in\mathfrak{p}\Rightarrow(a\in\mathfrak{p}\lor b\in\mathfrak{p})]$.
\end{definition}

\begin{theorem} \label{thm:prime-ideal-quotient-domain} \hlt{Prime Ideals and Quotient Domains}\\
Let $R$ be a commutative ring with identity and $\mathfrak{p}\trianglelefteq R$. Then $\mathfrak{p}$ is prime $\iff$ $R/\mathfrak{p}$ is an integral domain.
\begin{proof}
\begin{enumerate}[label=\roman*.]
\setlength{\itemsep}{0pt}
\item[$(\Rightarrow)$] Assume $\mathfrak{p}$ is prime. The ring $R/\mathfrak{p}$ is commutative with identity $1_R+\mathfrak{p}$ by Definition \ref{def:quotient-ring}.\\
The relation $1_R\notin\mathfrak{p}$ holds by $\mathfrak{p}\neq R$ and ideal absorption. The relation $1_R+\mathfrak{p}\neq\mathfrak{p}$ holds by additive-coset equality.\\
Fix $a,b\in R$ and assume $(a+\mathfrak{p})(b+\mathfrak{p})=\mathfrak{p}$. The relation $ab\in\mathfrak{p}$ holds by quotient multiplication and additive-coset equality.\\
The relation $a\in\mathfrak{p}$ or $b\in\mathfrak{p}$ holds by Definition \ref{def:prime-ideal}. The relation $a+\mathfrak{p}=\mathfrak{p}$ or $b+\mathfrak{p}=\mathfrak{p}$ holds by additive-coset equality.\\
The ring $R/\mathfrak{p}$ has no zero divisors by this zero-product test. It is an integral domain by Definition \ref{def:zero-divisor-integral-domain}.
\item[$(\Leftarrow)$] Assume $R/\mathfrak{p}$ is an integral domain. The relation $\mathfrak{p}\neq R$ holds by its distinct zero and identity cosets.\\
Fix $a,b\in R$ and assume $ab\in\mathfrak{p}$. The equality $(a+\mathfrak{p})(b+\mathfrak{p})=\mathfrak{p}$ holds by quotient multiplication and additive-coset equality.\\
The relation $a+\mathfrak{p}=\mathfrak{p}$ or $b+\mathfrak{p}=\mathfrak{p}$ holds by absence of zero divisors. Relation $a\in\mathfrak{p}$ or $b\in\mathfrak{p}$ holds by additive-coset equality.\\
The ideal $\mathfrak{p}$ is prime by Definition \ref{def:prime-ideal}.
\end{enumerate}
\hlt{Structural Commentary}: Quotient multiplication transfers the zero-divisor test to product membership in $\mathfrak{p}$.\\
\hlt{Key Insight}: A zero product of cosets is precisely a product in $\mathfrak{p}$.
\end{proof}
\end{theorem}

\begin{corollary} \label{cor:maximal-ideals-are-prime} \hlt{Maximal Ideals Are Prime}\\
Let $R$ be a commutative ring with identity. Every maximal ideal of $R$ is prime.
\begin{proof}
Let $\mathfrak{m}\trianglelefteq R$ be maximal. The quotient $R/\mathfrak{m}$ is a field by Theorem \ref{thm:quotient-by-maximal-ideal-is-field}.\\
Fix $x,y\in R/\mathfrak{m}$ with $xy=0$ and $x\neq0$. The equality $y=x^{-1}(xy)=0$ holds by field inversion.\\
The quotient $R/\mathfrak{m}$ is an integral domain by Definition \ref{def:zero-divisor-integral-domain}. The ideal $\mathfrak{m}$ is prime by Theorem \ref{thm:prime-ideal-quotient-domain}.\\
\hlt{Structural Commentary}: The quotient-field theorem feeds the quotient-domain criterion.\\
\hlt{Key Insight}: Invertibility excludes zero divisors in every field.
\end{proof}
\end{corollary}

\begin{definition} \label{def:ideal-radical} \hlt{Radicals and Radical Ideals}\\
Let $R$ be a commutative ring with identity and $\mathcal{I}\trianglelefteq R$.\\
The \hlt{radical} of $\mathcal{I}$ is $\operatorname{rad}(\mathcal{I})\defeq\sqrt{\mathcal{I}}\defeq\{r\in R:(\exists\,n\in\N)\,r^n\in\mathcal{I}\}$. The ideal $\mathcal{I}$ is \hlt{radical} if $\sqrt{\mathcal{I}}=\mathcal{I}$.
\end{definition}

\begin{proposition} \label{prop:ideal-radical-closure} \hlt{Radical Closure}\\
Let $R$ be a commutative ring with identity and $\mathcal{I}\trianglelefteq R$.
\begin{enumerate}[label=(\roman*)]
\setlength{\itemsep}{0pt}
\item $\sqrt{\mathcal{I}}\trianglelefteq R$.
\item $\mathcal{I}\subseteq\sqrt{\mathcal{I}}$.
\item The ideal $\sqrt{\mathcal{I}}$ is radical: $\operatorname{rad}(\operatorname{rad}(\mathcal{I}))=\operatorname{rad}(\mathcal{I})$.
\item If $\mathcal{J}\trianglelefteq R$ is radical and $\mathcal{I}\subseteq\mathcal{J}$, then $\sqrt{\mathcal{I}}\subseteq\mathcal{J}$.
\end{enumerate}
\begin{proof}
\begin{enumerate}[label=(\roman*)]
\setlength{\itemsep}{0pt}
\item The relation $0_R\in\sqrt{\mathcal{I}}$ holds by $0_R^1\in\mathcal{I}$. Fix $x,y\in\sqrt{\mathcal{I}}$ and choose $m,n\in\N$ with $x^m,y^n\in\mathcal{I}$.\\
For $0\leq k\leq m+n$, either $k\geq m$ or $m+n-k\geq n$; the relation $x^k(-y)^{m+n-k}\in\mathcal{I}$ holds by ideal absorption.\\
The relation $(x-y)^{m+n}\in\mathcal{I}$ holds by distributivity, ideal absorption, and additive closure; $x-y\in\sqrt{\mathcal{I}}$ by Definition \ref{def:ideal-radical}.\\
For $r\in R$, $(rx)^m=r^mx^m\in\mathcal{I}$ by commutativity and ideal absorption. The relation $rx=xr\in\sqrt{\mathcal{I}}$ holds by Definition \ref{def:ideal-radical}.\\
The conclusion $\sqrt{\mathcal{I}}\trianglelefteq R$ holds by Definition \ref{def:ideal-and-generated-ideal}.
\item Fix $x\in\mathcal{I}$; $x^1=x\in\mathcal{I}$. The relation $x\in\sqrt{\mathcal{I}}$ holds by Definition \ref{def:ideal-radical}. The inclusion $\mathcal{I}\subseteq\sqrt{\mathcal{I}}$ holds by this element test.
\item Fix $x\in\operatorname{rad}(\mathcal{I})$; $x^1\in\operatorname{rad}(\mathcal{I})$. The relation $x\in\operatorname{rad}(\operatorname{rad}(\mathcal{I}))$ holds by Definition \ref{def:ideal-radical}.\\
Inclusion $\operatorname{rad}(\mathcal{I})\subseteq\operatorname{rad}(\operatorname{rad}(\mathcal{I}))$ holds by this element test. If $x\in\operatorname{rad}(\operatorname{rad}(\mathcal{I}))$, choose $m\in\N$ with $x^m\in\operatorname{rad}(\mathcal{I})$.\\
Relation $(\exists\,n\in\N)\,x^{mn}=(x^m)^n\in\mathcal{I}$ holds by Definition \ref{def:ideal-radical}, exponent laws. Relation $x\in\operatorname{rad}(\mathcal{I})$ holds by Definition \ref{def:ideal-radical}.\\
Equality $\operatorname{rad}(\operatorname{rad}(\mathcal{I}))=\operatorname{rad}(\mathcal{I})$ holds by mutual inclusion. Ideal $\sqrt{\mathcal{I}}$ is radical by Definition \ref{def:ideal-radical}.
\item Let $\mathcal{J}\trianglelefteq R$ be radical with $\mathcal{I}\subseteq\mathcal{J}$. If $x\in\sqrt{\mathcal{I}}$, choose $n\in\N$ with $x^n\in\mathcal{I}\subseteq\mathcal{J}$.\\
The relation $x\in\sqrt{\mathcal{J}}=\mathcal{J}$ holds by Definition \ref{def:ideal-radical}. The inclusion $\sqrt{\mathcal{I}}\subseteq\mathcal{J}$ holds by this element test.
\end{enumerate}
\hlt{Structural Commentary}: Power witnesses establish ideal closure, idempotence, and minimality under radical containment.\\
\hlt{Key Insight}: A sufficiently large power absorbs every term in a binomial expansion.
\end{proof}
\end{proposition}
